\section{Probabilistic Evaluation and Conformal Upgrades}

\subsection{Why probabilistic output matters}

A point forecast tells a trader \emph{where} the price lands.
A probabilistic forecast tells them \emph{how wide to quote the spread}.
Risk managers and portfolio optimisers need quantiles, not just P50.

EU energy contracts increasingly require scenario generation.
If you only ship P50, you cannot do that.

\subsection{What we evaluate}

Three forecast objects:

\begin{itemize}
  \item \textbf{P50} (median): evaluated with MAE, rMAE.
  \item \textbf{P10/P90 band}: evaluated with pinball loss and empirical coverage.
  \item \textbf{Calibration}: does the 80\% band actually cover 80\% of outcomes?
\end{itemize}

\subsection{Pinball loss}

Also called quantile loss or tick loss.
For quantile $\tau \in (0,1)$ and forecast error $e = y - \hat{q}_\tau$:

\[
  L_\tau(e) =
  \begin{cases}
    \tau \cdot e      & \text{if } e \geq 0 \\
    (\tau - 1) \cdot e & \text{if } e < 0
  \end{cases}
\]

P10 ($\tau=0.1$) penalises over-forecasting 9$\times$ more than
under-forecasting. P90 is the mirror.
Sum the two: $L_{0.1} + L_{0.9}$. Lower is better.
Scale-independent comparison: divide by naive pinball.

In this project: LGBM pinball 0.124 vs naive 0.194 (36\% lower).

\subsection{Empirical coverage}

Count the fraction of hours where the actual price falls inside
$[\hat{q}_{0.1}, \hat{q}_{0.9}]$.
A well-calibrated model: coverage $\approx 80\%$.

Raw LGBM quantile regression: 51.4\% coverage on nominal 80\% band.
After conformal calibration: 78.7\%.
The gap shows raw quantile regression can be poorly calibrated.

\subsection{Conformal prediction recap}

\emph{Split conformal} (from note 16, summary):

\begin{enumerate}
  \item Calibration set: walk-forward out-of-sample errors from last 90 days.
  \item Compute nonconformity scores: $s_t = \max(\hat{q}_{0.1,t} - y_t,\ y_t - \hat{q}_{0.9,t})$.
  \item Find $\hat{q}_{0.9}(S)$ over all calibration scores.
  \item At test time: widen band by $\hat{q}_{0.9}(S)$ on each side.
\end{enumerate}

Coverage guarantee: $\Pr[y \in \text{band}] \geq 1 - \alpha = 0.80$.
No distributional assumption needed.

\subsection{Conformal Quantile Regression (CQR) — the upgrade}

Symmetric conformal widens both tails by the same amount.
CQR (Romano, Patterson \& Cand{\`e}s, 2019) is asymmetric:

\begin{align}
  s_t &= \max\!\bigl(\hat{q}_{0.1,t} - y_t,\ y_t - \hat{q}_{0.9,t}\bigr)
\end{align}

The key difference: CQR keeps the asymmetry learned by the quantile model.
If P90 is already well-calibrated but P10 undershoots,
CQR widens the lower tail more than the upper tail.
Symmetric CP cannot do this.

Gap in this project: symmetric CP corrects overall coverage from 51\% to 79\%.
It does not fix the asymmetric miscalibration that remains.
CQR is the next upgrade.

\subsection{Winkler score}

One number that combines sharpness and coverage:

\[
  W = (q_{0.9} - q_{0.1}) + \frac{2}{\alpha}
      \bigl[\max(q_{0.1}-y,0) + \max(y-q_{0.9},0)\bigr]
\]

Narrow band = low width term.
Missed coverage = high penalty term.
Lower is better. Useful for comparing conformal vs CQR vs other calibrators.

Limitation: if coverage drops below $1-\alpha$, the penalty
dominates and width differences become invisible.
Use Winkler only when comparing models with similar coverage.

\subsection{Interview lines}

\begin{itemize}
  \item ``Our raw quantile model achieved 51\% coverage on a nominal 80\% band.
        Conformal calibration fixed it to 79\% with no model retraining.''
  \item ``Pinball loss decomposes into two weighted MAE terms — one for
        over-forecasting, one for under-forecasting. The weights are $\tau$
        and $1-\tau$.''
  \item ``Symmetric conformal assumes both tails need the same correction.
        CQR is the right fix when one tail is already calibrated.''
  \item ``Winkler score combines sharpness and coverage in one number.
        It is the natural metric for comparing interval forecasts.''
\end{itemize}

\subsection{Practical checklist for shipping a probabilistic model}

\begin{enumerate}
  \item Plot empirical coverage by hour-of-day. Coverage should be flat.
        Spikes or dips reveal regime-specific miscalibration.
  \item Check coverage on high-price hours separately (e.g.\ $y > 150$ EUR/MWh).
        Most models under-cover extremes even with good overall coverage.
  \item Use rolling calibration (not a fixed window) so corrections adapt
        to regime shifts (e.g.\ winter gas crisis, summer solar oversupply).
  \item Log the conformal offset daily. A sudden jump indicates the model
        has entered a new regime.
  \item Never quote calibrated quantiles to a trading system without a
        stale-offset guard: if today's calibration data is missing,
        fall back to the last known-good offset, not to uncalibrated output.
\end{enumerate}
